\begin{abstract}
Ten instruments on four missions have mapped the Moon, but their products sit in the archive as raw per-image records and global maps with incompatible projections, resolutions and units, and labels are scarce. We present \method, a latent world model of the lunar surface, together with the co-registered corpus we build to train it. From these raw products we build a corpus of 10,629 footprints of 60\,km anchored to the GRAIL gravity grid, with 30 global map layers cut once per footprint at native resolution. Each footprint carries up to 500 Wide Angle Camera observations chosen for illumination diversity, each with its acquisition geometry, and 500\,m sub-tiles with 1\,m Narrow Angle Camera images and 3\,m stereo terrain; splits follow map zones and share no ground. \method treats a footprint as a persistent state seen through instruments. An encoder estimates the state from one or more partially masked observations, and a predictor, told which product to predict and the sun under which it was recorded, predicts the token grid that a moving-average encoder produces for that product. Every product is tokenized at its native resolution, so that a token covers the same ground cell in every product. Training this way has a failure mode the loss hides: a position code, in which tokens encode where they are rather than what the ground looks like, while the loss, within-image rank and VICReg penalties all improve. We explain why these signals are blind, show that two common regularizers push toward the shortcut, and give a cross-sample monitor that flagged every failure we analyzed. Frozen \method features detect craters on par with a masked-token lunar foundation model and gain 0.23 mAP$_{50}$ from a second input product, probes recover relief and albedo from a single image, and predicted shadows follow the sun the predictor is given (AUC 0.89 with the true sun, 0.59 with another).
\end{abstract}
